\documentclass[11pt,a4paper]{article}
\usepackage[margin=2.5cm]{geometry}
\usepackage{times}

\begin{document}

\thispagestyle{empty}

\begin{center}
{\Large \bf MNRAS --- Manuscript Submission}\\[8pt]
{\large Cover Letter}\\
\end{center}

\vspace{12pt}

\noindent\underline{Manuscript Details:}\\[4pt]
{\small
\textbf{Title:} A Redshift-Dependent Spectral Distortion in 260 Little Red Dots: Constraints on Compactness Evolution at Cosmic Dawn\\[3pt]
\textbf{Journal:} Monthly Notices of the Royal Astronomical Society (MNRAS)\\[3pt]
\textbf{Article Type:} Regular Article\\[3pt]
\textbf{Corresponding Author:} Xin Tan\\
Email: Summicron352@gmail.com\\
Affiliation: Institute for Advanced Information Physics (IAIP)
}

\vspace{16pt}

\noindent\underline{Dear Editor,}\\[6pt]

We are pleased to submit our manuscript entitled
``A Redshift-Dependent Spectral Distortion in 260 Little Red Dots:
Constraints on Compactness Evolution at Cosmic Dawn''
for consideration as a Regular Article in {\it Monthly Notices of the
Royal Astronomical Society}.

\vspace{6pt]

This work presents a comprehensive spectral energy distribution
(SED) analysis of the complete JWST/NIRSpec-confirmed sample of
260 Little Red Dots (LRDs; Kokorev et al.\ 2024). LRDs ---
compact, red sources at $z \sim 4$--$7$ exhibiting anomalous
V-shaped rest-frame optical continua --- have emerged as one of
the most puzzling populations discovered by {\it James Webb Space
Telescope}. Their spectral morphology defies straightforward
explanation within standard dust-attenuation or nebular continuum
frameworks.

\vspace{4pt]

We introduce a model-agnostic approach that treats a wavelength-
independent gravitational redshift component $z_{\rm dist}$ as a free
spectral distortion parameter alongside conventional dust extinction.
Fitting 260 LRD SEDs with both models yields three principal findings:

\vspace{4pt]

\noindent{\bf (1)} {\it Statistical preference for the Relic model.}
83.5\% of sources (217/260) show $\Delta\chi^2_\nu > -2$, indicating
that a non-dispersive redshift origin ($z_{\rm dist} = 0.1700 \pm 0.0015$
for the Strong Support subsample) provides a superior or equivalent
description compared to chromatic dust attenuation. The Support group
($\Delta\chi^2_\nu > 0$, 108 sources = 41.5\%) additionally exhibits
systematically reduced best-fit extinction ($A_V \to 0$), consistent
with a physical mechanism that replaces rather than supplements dust.

\vspace{4pt]

\noindent{\bf (2)} {\it Unified multi-wavelength resolution.} The
gravitational relic framework naturally explains four otherwise
independent observational puzzles simultaneously:
(a) V-shaped SEDs without requiring extreme $A_V = 2$--$4$;
(b) MIR silence despite high apparent obscuration (34/34 deep
MIRI non-detections + 0/260 Chandra detections);
(c) sub-virial emission-line widths ($\sim$3000\,km/s vs.\ theoretical
$\sim$\!100,000\,km/s under virial equilibrium); and
(d) systematically low stellar masses from photometric fitting,
consistent with flux dilution by $1+z_{\rm dist}$.

\vspace{4pt]

\noindent{\bf (3)} {\it Statistical impossibility of independent
coincidence.} Under the standard dust model, six independent rare
conditions must coincide to reproduce the observed pattern, yielding a
joint probability $P_{\rm joint} \approx 1.8 \times 10^{-6}$ per source.
This predicts $\sim$\!0.03 sources out of 260, whereas 217 (83.5\%)
actually exhibit the full set of anomalous features --- a binomial
deviation exceeding $100\sigma$. This over-constraint argument rules out
chance coincidence and points toward a single unified physical origin.

\vspace{6pt]

The hypothesis is falsifiable through two near-term observations:
(i) spatially resolved stellar velocity dispersion measurements via
JWST/NIRSpec IFU (prediction: $\sigma_* > 300$\,km/s, inconsistent
with normal star-forming galaxies at this epoch); and (ii) detection
of anomalous strong gravitational lensing signals around individual
LRDs (prediction: central convergence $\kappa_0 \gg 1$).

\vspace{6pt]

We believe this work will be of broad interest to the MNRAS readership
engaged with JWST extragalactic science, high-redshift galaxy physics,
and alternative interpretations of early-universe phenomena. The analysis
is fully reproducible, and we commit to making all fitting codes publicly
available upon acceptance.

\vspace{12pt]

\noindent We suggest the following potential reviewers:

\vspace{4pt]

{\small
\begin{itemize}
\item[(i)] Prof.\ Pablo G.\ P\'{e}rez-Gonz\'{a}lez (UCM/CAB, Spain) --- expert in NIRCam-selected galaxy populations and SED modelling;
\item[(ii)] Prof.\ Roberto Maiolino (University of Cambridge, UK) --- leading authority on LRD physics and JWST spectroscopy;
\item[(iii)] Prof.\ Jeyhan Kartaltepe (Rochester Institute of Technology, USA) --- extensive experience with AGN-galaxy co-evolution and JWST surveys;
\item[(iv)] Prof.\ Hannah \"Ubler (ETH Zurich, Switzerland) --- specialist in emission-line kinematics of high-redshift compact galaxies.
\end{itemize}
}

\vspace{14pt]

This manuscript has not been published and is not under consideration
for publication elsewhere. All authors have read and approved the submission.

\vspace{18pt]

\noindent Yours sincerely,

\vspace{8pt]

\noindent\underline{Xin Tan}\\[3pt]
{\small Independent Researcher}\\[3pt]
{\small Institute for Advanced Information Physics (IAIP)}\\[3pt]
{\small Email: Summicron352@gmail.com}

\end{document}
